\subsection{The remaining star-shaped support reductions}

\begin{lemma}[Proved: coning the source of a binary bridge]
\label{CP2:A-star-coned-binary}
Let $S$ be a nonempty left-$N$, right-$M$ biset with at most
two states, for finite monoids $M,N$. Put $E=B(N,S,M)$.
Every rational degree-two product on
$\operatorname{hocofib}(E\to BN)$ is zero.
Consequently the same is true for a source star if one stalk
is a singleton with rationally acyclic leaf corner, and the
other stalk has at most two states.
\end{lemma}
\begin{proof}
A single state gives a section of $E\to BN$. Suppose there
are two states and inspect the right $M$ transformation image.
If it contains a constant map with value $p$, commutation
forces every $N$ actor to fix $p$. Using this state and the
identity $M$ string gives a section $BN\to E$. If the image
contains no constant and contains the transposition, every
commuting $N$ actor is a permutation. In this case
$E\to BN\times BM$ is the finite covering associated to the
permutation action on $S$. Transfer, followed by the injection
from the $BN$ factor, makes $H^2(BN;\mathbb Q)\to H^2(E;\mathbb Q)$
injective. In either case the cofiber exact sequence makes
every second class vanish on $BN$. The suspension quotient
and relative-cup argument of Lemma~\ref{CP2:A-two-block} kill
all its degree-two products.

The remaining right image is the identity alone. Then
$E=BM\times Z$, where $Z=B(N,S)$ is the ordinary action space.
Projection to $Z$ induces a map of cofibers to the ordinary
two-state $N$-act cone $Y$. Since $H^1(BM;\mathbb Q)=0$,
it induces an isomorphism on $H^1$ of the bridge spaces.
The two kernels from $H^2(BN;\mathbb Q)$ agree: projection
$BM\times Z\to Z$ has a section. The cofiber exact sequences
therefore give an isomorphism $H^2(Y;\mathbb Q)\to H^2(X;\mathbb Q)$.
The affine two-state lemma, Lemma~\ref{CP2:A-high-affine},
kills products on $Y$, and naturality kills them on $X$.

For the stated star, its singleton edge is $BM*BN_0$ and
is rationally acyclic when $BN_0$ is. Collapse that edge.
The quotient map is a rational cohomology-ring isomorphism,
and the resulting space is exactly the cofiber just treated.
\end{proof}

\begin{corollary}[Proved: a center of order at most four]
\label{CP2:A-star-small-center}
A three-object star with at most seventeen arrows and center
of order at most four cannot have type $\CP^2$.
\end{corollary}
\begin{proof}
By Proposition~\ref{CP2:A-star-large-leaf}, assume both
leaves have order at most seven. The small-center argument
in that proposition applies without its large-leaf
assumption. If the center classifying space is contractible,
the star is the wedge of its two full edges. A projective
plane target would force one of these edges to be an actual
projective-plane model, with at most fifteen arrows,
contradicting the two-object baseline. If the center has
nontrivial minimal group, that proof shows that its minimal
ideal is a group and both trivial rational modules are
projective. The actual orbit quotient at the center gives
a rationally equivalent star with trivial center. Its two
ordinary act-cone wedge summands have bounded rational
homology and corner orders at most seven. The bounded-nine
act theorem kills their degree-two products, and hence
those of the whole wedge.
\end{proof}

\begin{lemma}[Proved: two small leaves and star retracts]
\label{CP2:A-star-leaf-reductions}
The following reductions hold for a source star.
\begin{enumerate}
\item A leaf corner whose trivial right rational module is
projective can be replaced by the trivial corner, with its
stalk replaced by the actual left orbit set. This gives a
rational homology and cohomology-ring isomorphism.
\item Suppose a leaf corner has order at most four and its
stalk has at most two states. Either the center fixes a
state of that stalk, or that leaf can be replaced as in (1)
by a trivial corner, with at most two states.
\item If both leaves are trivial, identifying their two
vertices gives the ordinary center-act cone on the disjoint
union of the stalks. Its classifying space is homotopy
equivalent to the original star space wedged with a circle.
\item Suppose a star of rational homology and cohomology
ring $\CP^2$ retracts onto a subcategory, and its quotient
by that subcategory is a suspension. Then the retract also
has the rational homology and cohomology ring of $\CP^2$.
This applies both to an equivariant retraction of one
stalk, leaving the other stalk unchanged, and to merging
one trivial leaf into the other along an equivariant map
of their center acts.
\end{enumerate}
\end{lemma}
\begin{proof}
Part (1) is the actual orbit quotient of
Lemma~\ref{CP2:A-high-orbits}, applied at that leaf. To prove
(2), only the nonprojective-right case remains. The small
minimal-ideal description shows that the minimal group is
trivial and the ideal is a one-row rectangular band. Indeed,
two rows and two columns already require five elements,
while one column makes the trivial right module projective;
a nontrivial minimal group with two columns requires at
least five elements. The left actors of this one-row band
are retractions with a common image $U$. On a set of at
most two states, either $|U|=1$ or all these retractions
are identities. In the first case the sole point of $U$
is fixed by the commuting center action. In the second,
every actor of the whole leaf monoid is the identity:
left multiplication fixes every element of a one-row band,
so its action fixes the image $U=S$. Its classifying space
is contractible by the small minimal-ideal lemma. The
quotient replacing it by a trivial corner is then a
homotopy equivalence in the double bar model.

For (3), the quotient identifying the leaf vertices is
the nerve realization of the ordinary act cone on the
disjoint union. Identifying two distinct points in a
connected CW complex has homotopy type $X\vee S^1$.
Equivalently, first attach an interval between the points
and then collapse that contractible interval.

For (4), if the second cohomology of the retract were zero,
all second classes would lift from the suspension quotient,
and their products would vanish. Thus the retract has
nonzero second cohomology. Its homology is a direct summand
of that of the original space, so it has one-dimensional
second homology and no other positive homology except
possibly in degree four. Pullback in second cohomology is
an isomorphism. The nonzero square forces a nonzero fourth
group in the retract, hence its entire rational target
pattern and ring. For a retraction of one stalk the quotient
is the suspension of the quotient of the corresponding
two-sided action spaces; all corner spaces and the other
edge are unchanged. For merging two trivial leaves the
retract is the retained edge, and the quotient is the
suspension of the removed stalk's ordinary action space.
These are literal quotients of the mapping-cylinder models.
\end{proof}

\begin{proposition}[Proved: a large center in a small star]
\label{CP2:A-star-large-center}
A connected three-object category with star-shaped acyclic
support, at most seventeen arrows, and center corner of
order at least eight cannot have type $\CP^2$.
\end{proposition}
\begin{proof}
Take opposition so the center is a source. Write its
order as $a$, its leaf orders as $b,c$, and its stalk sizes
as $s,t$. Thus $a+b+c+s+t\le17$. Large leaves are already
excluded by Proposition~\ref{CP2:A-star-large-leaf}; assume
$b,c\le7$.

Suppose first $a\ge10$. Both leaves have order at most
four, hence are rationally acyclic. A singleton stalk gives
a center-fixed state, and Corollary~\ref{CP2:A-fixed-arrow}
in its dual form kills every degree-two product. We may
assume $s,t\ge2$. The only possibilities are $a=11$ with
$b=c=1$ and $s=t=2$, or $a=10$ with leaf orders $(1,1)$
and stalk sizes at most $(2,3)$, or leaf orders $(1,2)$
and stalk sizes $(2,2)$, up to interchange. Orbit comparison
at the two-element leaf reduces the latter case to trivial
leaves. If an orbit set becomes a singleton, the same fixed
state argument applies. Otherwise the reduced star has the
rational homology and ring of $\CP^2$ and two trivial leaves,
one on two states and the other on at most three.

Merge the leaves as in Lemma~\ref{CP2:A-star-leaf-reductions}(3).
The resulting ordinary act cone has bounded rational
homology and the same positive-degree products in degrees
at least two. Relative cups require $H^2(BM;\mathbb Q)\ne0$;
the two-row obstruction requires at least three minimal
rows, and there are at least two columns. The cardinality
budget makes the minimal group trivial. If there are at
least three columns, the ideal is $3\times3$ and there
are at most two outside elements. Lemma~\ref{CP2:A-two-outside-acts}
then makes fourth homology of this bounded cone zero,
a contradiction.

Thus the minimal band has two columns. A band-fixed state
is center-fixed: if $xk=x$ for every $k$ in the ideal, then
$xm=xkm=x$ for every $m$. This again kills all products.
In its absence every core orbit has two points. On every
such orbit the action of $M$ is the same labelled column
action, because right multiplication sends $e_{ij}$ to
$e_{i,\rho_m(j)}$ independently of $i$. The two-state
stalk is one such orbit and therefore embeds equivariantly
into a core orbit of the other stalk. Merge its leaf into
the other leaf along this embedding. Part (4) of the
preceding lemma makes the retained ordinary act cone a
rational $\CP^2$ target. If $a=11$, it has two states and
the affine lemma excludes it. If $a=10$, it has at most
three states and Proposition~\ref{CP2:A-ten-six-state-rational}
excludes it. This finishes $a\ge10$.

Now let $a=8$ or $9$. If a leaf has order at least five,
the other has order at most two and $s+t\le3$. If the
large leaf's stalk is a singleton, that edge is a join
with zero second cohomology. It kills restrictions to
the center and large leaf; the remaining leaf is rationally
acyclic. The two-block relative-cup lemma applies. If the
small leaf has the singleton stalk, the other has at most
two states, and Lemma~\ref{CP2:A-star-coned-binary} applies.

We are left with $b,c\le4$. A singleton stalk again gives
a fixed state and zero products. Assume $s,t\ge2$. If
both leaf orders are at most two, orbit comparison makes
them trivial and Lemma~\ref{CP2:A-star-leaf-reductions}(3)
reduces to the bounded-nine act theorem. Every leaf on
two states can be treated in the same way by part (2),
since its fixed-state alternative kills products. The
only budget case not yet covered has
\[
 (a,b,c,s,t)=(8,3,1,3,2)
\]
up to interchange of the leaves. If the three-element
leaf has projective trivial right module, its orbit
quotient once again gives the bounded-nine act case.
Otherwise it is $R_0=\{1,p,q\}$, where $uv=v$ on the
two nonidentity elements.

Its two actors $P,Q$ on the three-state stalk are
retractions with the same image $U$. If $P=Q$, their
common retraction is equivariant for both corner actions.
If $|U|=1$, its point is center-fixed. Otherwise retract
this stalk to $U$. Part (4) transfers the rational target
to the retracted star. The $R_0$ action on $U$ is trivial
and $BR_0$ is contractible, so quotienting that leaf to
the trivial corner gives a rational equivalence (indeed,
a homotopy equivalence). The bounded-nine act theorem
after merging leaves again excludes the result.

If $P\ne Q$, necessarily $|U|=2$. The natural graph model
of the $R_0$ action category in
Lemma~\ref{CP2:A-high-P5-graph} has the two vertices $U$
and one edge from $P(x)$ to $Q(x)$ for the third state
$x$. It is an interval; its construction commutes with
the center action. Its homotopy-orbit space over $M$
therefore maps by a homotopy equivalence to $BM$.
Quotient the entire stalk to one state and the leaf to
the trivial monoid. The map on bridge spaces is this
equivalence, on the leaf spaces it is $BR_0\to *$, and
on the center it is the identity. Homotopy-pushout
invariance gives a homotopy equivalence of stars. Its
singleton stalk now supplies a center-fixed state, so
its degree-two products vanish. All cases are excluded.
\end{proof}
